% Options for packages loaded elsewhere
% Options for packages loaded elsewhere
\PassOptionsToPackage{unicode}{hyperref}
\PassOptionsToPackage{hyphens}{url}
\PassOptionsToPackage{dvipsnames,svgnames,x11names}{xcolor}
%
\documentclass[
  letterpaper,
  DIV=11,
  numbers=noendperiod]{scrartcl}
\usepackage{xcolor}
\usepackage{amsmath,amssymb}
\setcounter{secnumdepth}{-\maxdimen} % remove section numbering
\usepackage{iftex}
\ifPDFTeX
  \usepackage[T1]{fontenc}
  \usepackage[utf8]{inputenc}
  \usepackage{textcomp} % provide euro and other symbols
\else % if luatex or xetex
  \usepackage{unicode-math} % this also loads fontspec
  \defaultfontfeatures{Scale=MatchLowercase}
  \defaultfontfeatures[\rmfamily]{Ligatures=TeX,Scale=1}
\fi
\usepackage{lmodern}
\ifPDFTeX\else
  % xetex/luatex font selection
\fi
% Use upquote if available, for straight quotes in verbatim environments
\IfFileExists{upquote.sty}{\usepackage{upquote}}{}
\IfFileExists{microtype.sty}{% use microtype if available
  \usepackage[]{microtype}
  \UseMicrotypeSet[protrusion]{basicmath} % disable protrusion for tt fonts
}{}
\makeatletter
\@ifundefined{KOMAClassName}{% if non-KOMA class
  \IfFileExists{parskip.sty}{%
    \usepackage{parskip}
  }{% else
    \setlength{\parindent}{0pt}
    \setlength{\parskip}{6pt plus 2pt minus 1pt}}
}{% if KOMA class
  \KOMAoptions{parskip=half}}
\makeatother
% Make \paragraph and \subparagraph free-standing
\makeatletter
\ifx\paragraph\undefined\else
  \let\oldparagraph\paragraph
  \renewcommand{\paragraph}{
    \@ifstar
      \xxxParagraphStar
      \xxxParagraphNoStar
  }
  \newcommand{\xxxParagraphStar}[1]{\oldparagraph*{#1}\mbox{}}
  \newcommand{\xxxParagraphNoStar}[1]{\oldparagraph{#1}\mbox{}}
\fi
\ifx\subparagraph\undefined\else
  \let\oldsubparagraph\subparagraph
  \renewcommand{\subparagraph}{
    \@ifstar
      \xxxSubParagraphStar
      \xxxSubParagraphNoStar
  }
  \newcommand{\xxxSubParagraphStar}[1]{\oldsubparagraph*{#1}\mbox{}}
  \newcommand{\xxxSubParagraphNoStar}[1]{\oldsubparagraph{#1}\mbox{}}
\fi
\makeatother


\usepackage{longtable,booktabs,array}
\usepackage{calc} % for calculating minipage widths
% Correct order of tables after \paragraph or \subparagraph
\usepackage{etoolbox}
\makeatletter
\patchcmd\longtable{\par}{\if@noskipsec\mbox{}\fi\par}{}{}
\makeatother
% Allow footnotes in longtable head/foot
\IfFileExists{footnotehyper.sty}{\usepackage{footnotehyper}}{\usepackage{footnote}}
\makesavenoteenv{longtable}
\usepackage{graphicx}
\makeatletter
\newsavebox\pandoc@box
\newcommand*\pandocbounded[1]{% scales image to fit in text height/width
  \sbox\pandoc@box{#1}%
  \Gscale@div\@tempa{\textheight}{\dimexpr\ht\pandoc@box+\dp\pandoc@box\relax}%
  \Gscale@div\@tempb{\linewidth}{\wd\pandoc@box}%
  \ifdim\@tempb\p@<\@tempa\p@\let\@tempa\@tempb\fi% select the smaller of both
  \ifdim\@tempa\p@<\p@\scalebox{\@tempa}{\usebox\pandoc@box}%
  \else\usebox{\pandoc@box}%
  \fi%
}
% Set default figure placement to htbp
\def\fps@figure{htbp}
\makeatother





\setlength{\emergencystretch}{3em} % prevent overfull lines

\providecommand{\tightlist}{%
  \setlength{\itemsep}{0pt}\setlength{\parskip}{0pt}}



 


\KOMAoption{captions}{tableheading}
\makeatletter
\@ifpackageloaded{caption}{}{\usepackage{caption}}
\AtBeginDocument{%
\ifdefined\contentsname
  \renewcommand*\contentsname{Table of contents}
\else
  \newcommand\contentsname{Table of contents}
\fi
\ifdefined\listfigurename
  \renewcommand*\listfigurename{List of Figures}
\else
  \newcommand\listfigurename{List of Figures}
\fi
\ifdefined\listtablename
  \renewcommand*\listtablename{List of Tables}
\else
  \newcommand\listtablename{List of Tables}
\fi
\ifdefined\figurename
  \renewcommand*\figurename{Figure}
\else
  \newcommand\figurename{Figure}
\fi
\ifdefined\tablename
  \renewcommand*\tablename{Table}
\else
  \newcommand\tablename{Table}
\fi
}
\@ifpackageloaded{float}{}{\usepackage{float}}
\floatstyle{ruled}
\@ifundefined{c@chapter}{\newfloat{codelisting}{h}{lop}}{\newfloat{codelisting}{h}{lop}[chapter]}
\floatname{codelisting}{Listing}
\newcommand*\listoflistings{\listof{codelisting}{List of Listings}}
\makeatother
\makeatletter
\makeatother
\makeatletter
\@ifpackageloaded{caption}{}{\usepackage{caption}}
\@ifpackageloaded{subcaption}{}{\usepackage{subcaption}}
\makeatother
\usepackage{bookmark}
\IfFileExists{xurl.sty}{\usepackage{xurl}}{} % add URL line breaks if available
\urlstyle{same}
\hypersetup{
  pdftitle={Midterm Project - GitHub},
  colorlinks=true,
  linkcolor={blue},
  filecolor={Maroon},
  citecolor={Blue},
  urlcolor={Blue},
  pdfcreator={LaTeX via pandoc}}


\title{Midterm Project - GitHub}
\usepackage{etoolbox}
\makeatletter
\providecommand{\subtitle}[1]{% add subtitle to \maketitle
  \apptocmd{\@title}{\par {\large #1 \par}}{}{}
}
\makeatother
\subtitle{JSC370 (Winter 2026)}
\author{}
\date{}
\begin{document}
\maketitle


\subsubsection{\texorpdfstring{Github
\href{link}{https://github.com/Richard098789/JSC370/tree/main/final\_project}}{Github https://github.com/Richard098789/JSC370/tree/main/final\_project}}\label{github-httpsgithub.comrichard098789jsc370treemainfinal_project}

\section{Introduction}\label{introduction}

Online job postings provide valuable insights into labour market demand
and skills. As a data science student currently looking for a job, I am
interested in \textbf{how do role type, location, and job
characteristics relate to salary levels in Canadian data-related job
postings? Does salary increase over time?}.

In this project, job posting data were collected using the Adzuna Job
Search API, which collects job listings from multiple job boards and
company career pages. Each job posting contains information such as the
job title, company name, location, salary range, contract type, and
description. These attributes make the dataset suitable for exploratory
data analysis, as it includes a mix of numerical, binary and categorical
variables.

The goal of this project is to explore patterns in salary levels across
different job characteristics. In particular, the analysis focuses on
comparing salaries across major data-related roles such as data
scientists, data engineers, and data analysts, as well as examining
geographic differences across Canadian regions. Additional exploratory
analysis investigates whether the of job posting date is associated with
variations in salary levels.

\section{Method}\label{method}

\subsubsection{1. Data Collection}\label{data-collection}

\begin{itemize}
\item
  Use Adzuna Job Search API to collect job postings data. Restrict the
  search range to Canada and the job title to include ``data''. Convert
  the data into a pandas dataframe. Fetch the first num\_pages pages of
  data, with 50 jobs per page.
\item
  Collected 5000 job postings, with features: title, company, location,
  salary\_min, salary\_max, contract\_type, category, created.
\end{itemize}

\subsubsection{\texorpdfstring{2. Load and clean (\texttt{data.csv} →
\texttt{df\_clean} →
\texttt{df\_feat})}{2. Load and clean (data.csv → df\_clean → df\_feat)}}\label{load-and-clean-data.csv-df_clean-df_feat}

Saved API pulls live in \texttt{data.csv}. The pipeline: (1) load and
de-duplicate; (2) keep rows with observed \texttt{salary\_min}, engineer
text and geography features; (3) attach CPI and restrict to postings on
or after \textbf{2025-04-01} for modeling and figures.

\begin{longtable}[]{@{}lll@{}}
\toprule\noalign{}
& stage & n \\
\midrule\noalign{}
\endhead
\bottomrule\noalign{}
\endlastfoot
0 & rows in data.csv & 5000 \\
1 & unique id (if present) & 5000 \\
\end{longtable}

\begin{longtable}[]{@{}ll@{}}
\toprule\noalign{}
& pct\_missing \\
\midrule\noalign{}
\endhead
\bottomrule\noalign{}
\endlastfoot
salary\_max & 76.0 \\
salary\_min & 73.9 \\
contract\_type & 72.7 \\
company & 0.0 \\
title & 0.0 \\
id & 0.0 \\
location\_detail & 0.0 \\
location & 0.0 \\
category & 0.0 \\
created & 0.0 \\
date & 0.0 \\
\end{longtable}

\begin{longtable}[]{@{}ll@{}}
\toprule\noalign{}
& n \\
\midrule\noalign{}
\endhead
\bottomrule\noalign{}
\endlastfoot
rows in data.csv & 5000 \\
unique id (if present) & 5000 \\
after dropna(salary\_min) & 1307 \\
\end{longtable}

\begin{longtable}[]{@{}llllll@{}}
\toprule\noalign{}
& title & title\_clean & role & province & salary\_mid \\
\midrule\noalign{}
\endhead
\bottomrule\noalign{}
\endlastfoot
23 & Data Engineer & data engineer & Data Engineer & Québec &
120000.0 \\
25 & Data Architect & data architect & Data Engineer &
Canada-unspecified & 105000.0 \\
27 & Data Steward & data steward & Other & Ontario & 124800.0 \\
29 & Data Analyst & data analyst & Data Analyst & Ontario & 125000.0 \\
32 & Data Analyst & data analyst & Data Analyst & British Columbia &
88832.0 \\
\end{longtable}

\begin{longtable}[]{@{}lllll@{}}
\toprule\noalign{}
& location\_detail & location & cpi & salary\_mid \\
\midrule\noalign{}
\endhead
\bottomrule\noalign{}
\endlastfoot
23 & Montréal & Québec & 167.0 & 120000.0 \\
25 & Canada & Canada & 167.4 & 105000.0 \\
27 & Toronto & Ontario & 169.7 & 124800.0 \\
29 & Toronto & Ontario & 169.7 & 125000.0 \\
32 & Greater Vancouver & British Columbia & 164.7 & 88832.0 \\
\end{longtable}

\section{Preliminary Results}\label{preliminary-results}

\subsubsection{0. Basic Summary
Statistics}\label{basic-summary-statistics}

\begin{figure}

\centering{

\centering{

\begin{verbatim}
Unable to display output for mime type(s): text/html
\end{verbatim}

}

\subcaption{\label{fig-sample-summary-1}Analytic sample of Canadian
data-related job postings (restricted date window; salary midpoints
annualized per cleaning rules).}

\centering{

\begin{verbatim}
Unable to display output for mime type(s): text/html
\end{verbatim}

}

\subcaption{\label{fig-sample-summary-2}}

}

\caption{\label{fig-sample-summary}}

\end{figure}%

The analytic sample summarized in the Plotly table above is the basis
for the tables and figures in this section unless otherwise noted (post
date filter and salary rules as under \emph{Methods}).

\subsubsection{1. Salary Distribution}\label{salary-distribution}

\begin{verbatim}
count      1285.000000
mean     104200.538911
std       52191.831951
min       20800.000000
25%       67500.000000
50%      105000.000000
75%      130000.000000
max      425000.000000
Name: salary_mid, dtype: float64
\end{verbatim}

Visualize the distribution of the salary midpoints.

\begin{verbatim}
Text(0, 0.5, 'Number of Job Postings')
\end{verbatim}

\pandocbounded{\includegraphics[keepaspectratio]{midterm_files/figure-pdf/cell-9-output-2.pdf}}

Figure shows the distribution of salary midpoint values. Most job
postings fall within a middle salary range, while a smaller number of
postings list substantially higher salaries.

\subsubsection{2. Salary Distribution by Job
Role}\label{salary-distribution-by-job-role}

\begin{longtable}[]{@{}lllll@{}}
\toprule\noalign{}
& count & mean & median & std \\
role & & & & \\
\midrule\noalign{}
\endhead
\bottomrule\noalign{}
\endlastfoot
Data Analyst & 247 & 99087.0 & 92900.0 & 34982.0 \\
Data Engineer & 242 & 133046.0 & 120000.0 & 49524.0 \\
Data Entry & 191 & 35745.0 & 29900.0 & 16402.0 \\
Data Scientist & 79 & 129945.0 & 125000.0 & 39352.0 \\
Developer & 61 & 124646.0 & 120000.0 & 38223.0 \\
Manager & 237 & 132881.0 & 126375.0 & 43916.0 \\
Other & 228 & 92266.0 & 80403.0 & 46110.0 \\
\end{longtable}

\begin{verbatim}
Text(0, 0.5, 'Salary Midpoint')
\end{verbatim}

\pandocbounded{\includegraphics[keepaspectratio]{midterm_files/figure-pdf/cell-11-output-2.pdf}}

Figure compares salary distributions across job roles. Data Engineer has
highest salary, then Data Scientist roles, with Data Analyst roles
having the lowest salary. Data Enguneer also has a high variance in
salary.

\subsubsection{3. Number of Job Postings by Job
Role}\label{number-of-job-postings-by-job-role}

\begin{verbatim}
Text(0, 0.5, 'Number of Job Postings')
\end{verbatim}

\pandocbounded{\includegraphics[keepaspectratio]{midterm_files/figure-pdf/cell-12-output-2.pdf}}

Figure shows the number of job postings by job role. Data Engineer has
the most job postings, followed by Data Analyst, and Data Scientist has
the least job postings.

\subsubsection{4. Salary Distribution by Contract
Type}\label{salary-distribution-by-contract-type}

\begin{verbatim}
Text(0, 0.5, 'Salary Midpoint')
\end{verbatim}

\pandocbounded{\includegraphics[keepaspectratio]{midterm_files/figure-pdf/cell-13-output-2.pdf}}

Figure shows the salary distribution by contract type. Permanent
contracts have around same salary compared to contract jobs.
(\textasciitilde100000)

\subsubsection{5. Salary Distribution by
Province}\label{salary-distribution-by-province}

\begin{longtable}[]{@{}llllll@{}}
\toprule\noalign{}
& province & count & mean & median & std \\
\midrule\noalign{}
\endhead
\bottomrule\noalign{}
\endlastfoot
0 & Alberta & 64 & 93930.0 & 88780.0 & 53595.0 \\
1 & British Columbia & 152 & 93302.0 & 90000.0 & 54171.0 \\
2 & Canada-unspecified & 275 & 129297.0 & 122500.0 & 54485.0 \\
3 & Ontario & 668 & 103477.0 & 107500.0 & 46979.0 \\
4 & Other Provinces & 57 & 71770.0 & 84435.0 & 36587.0 \\
5 & Québec & 69 & 71512.0 & 48880.0 & 51602.0 \\
\end{longtable}

\begin{verbatim}
([0, 1, 2, 3, 4, 5],
 [Text(0, 0, 'Québec'),
  Text(1, 0, 'Canada-unspecified'),
  Text(2, 0, 'Ontario'),
  Text(3, 0, 'British Columbia'),
  Text(4, 0, 'Alberta'),
  Text(5, 0, 'Other Provinces')])
\end{verbatim}

\pandocbounded{\includegraphics[keepaspectratio]{midterm_files/figure-pdf/cell-15-output-2.pdf}}

Figure shows the salary distribution by province. Jobs with region not
posted in Canada has the highest salary, followed by Ontario, Quebec,
and Manitoba, Nova Scotia, New Brunswick have lower salary.

\subsubsection{6. Average salary by posting date and job
role}\label{average-salary-by-posting-date-and-job-role}

\begin{verbatim}
(array([20209., 20270., 20332., 20393., 20454., 20513., 20574.]),
 [Text(20209.0, 0, '2025-05'),
  Text(20270.0, 0, '2025-07'),
  Text(20332.0, 0, '2025-09'),
  Text(20393.0, 0, '2025-11'),
  Text(20454.0, 0, '2026-01'),
  Text(20513.0, 0, '2026-03'),
  Text(20574.0, 0, '2026-05')])
\end{verbatim}

\pandocbounded{\includegraphics[keepaspectratio]{midterm_files/figure-pdf/cell-16-output-2.pdf}}

\begin{verbatim}
(array([20209., 20270., 20332., 20393., 20454., 20513., 20574.]),
 [Text(20209.0, 0, '2025-05'),
  Text(20270.0, 0, '2025-07'),
  Text(20332.0, 0, '2025-09'),
  Text(20393.0, 0, '2025-11'),
  Text(20454.0, 0, '2026-01'),
  Text(20513.0, 0, '2026-03'),
  Text(20574.0, 0, '2026-05')])
\end{verbatim}

\pandocbounded{\includegraphics[keepaspectratio]{midterm_files/figure-pdf/cell-17-output-2.pdf}}

Figure shows a salary trend by posting date and job role. All 3 roles
have similar salary trend, with a slight increase in salary over time.

\section{Summary}\label{summary}

After cleaning and processing the data, the analysis would focus on
several key features including job role, location, contract type and
created date.

The exploratory analysis revealed several preliminary patterns in the
dataset. First, salary distributions differ across job roles. Positions
categorized as Data Engineer and Data Scientist generally have higher
salary ranges compared to Data Analyst roles.

Second, the geographic distribution of job postings is uneven across
Canada. A large proportion of postings in the dataset are located in
Ontario, while smaller numbers appear in provinces such as British
Columbia, Quebec, and Alberta. Salary distributions also vary across
provinces, although the number of observations in some regions is
relatively small. These geographic differences suggest that regional job
markets may influence salary levels.

Third, the temporal trend analysis shows that average salary have been
increasing over time since 2024, with a slight decrease in 2025. Since
most job fetched are from recent time, if more job postings are
collected from older time, we can further investigate the trend.

\section{Plan}\label{plan}

First, I wll expand the data size by fetching more job postings using
Adzuna API to include more job postings from older time. This would help
learn more about the trend of salary over time.

Next I will use ML methods to further explore the relationship between
salary and other features. As an initial baseline, a linear regression
model will be fitted to estimate salary based on features such as job
role, geographic location, contract type, and posting date. To capture
potential nonlinear relationships, models such as GAM will also be used.
Finally, ensemble learning methods including random forests and gradient
boosting will be used to improve predictive performance and identify the
most important factors influencing salary.

I would focus on exploring role type, and whether different months of a
year have different salary levels. I would also focus on exploring
region and compare them with actual Canadian geographic boundaries to
better understand regional job market patterns.

The result of the model will be added to interactive website, where user
can pick the features they are interested in and see the predicted
salary.

\subsubsection{Predictor correlation (multicollinearity
screening)}\label{predictor-correlation-multicollinearity-screening}

Pearson correlations among the \textbf{training-set} modeling features
flag strong linear redundancy. \textbf{Random forest and XGBoost}
tolerate correlated predictors reasonably well for prediction;
\textbf{GAM} smooths on both \textbf{\texttt{year}} and
\textbf{\texttt{month}} partly overlap in what they explain---\textbf{do
not read them as independent additive ``year effect'' plus ``month
effect''} when \textbf{\texttt{year} and \texttt{month} are highly
correlated} (often (\textbar r\textbar{}\gtrsim 0.85) in a bounded
posting window because both trace \textbf{calendar time}). A tighter
design would use a \textbf{single time index} (e.g.~months or days since
start of the analytic window); here we \textbf{retain both} and note the
limitation. Strong (\textbar r\textbar) among other pairs appear in the
sorted list below ((\geq 0.65) rule-of-thumb rows).

\pandocbounded{\includegraphics[keepaspectratio]{midterm_files/figure-pdf/cell-19-output-1.pdf}}

\begin{longtable}[]{@{}llll@{}}
\toprule\noalign{}
& feat\_a & feat\_b & pearson\_r \\
\midrule\noalign{}
\endhead
\bottomrule\noalign{}
\endlastfoot
0 & year & month & -0.918519 \\
1 & is\_major\_city & missing\_province & -0.488867 \\
2 & role\_encoded & year & 0.436080 \\
3 & role\_encoded & is\_permanent & -0.417748 \\
4 & is\_permanent & month & 0.413676 \\
5 & role\_encoded & month & -0.381460 \\
6 & is\_permanent & year & -0.374935 \\
7 & is\_major\_city & cpi & 0.302729 \\
8 & role\_encoded & is\_senior & 0.285174 \\
9 & role\_encoded & is\_major\_city & 0.209868 \\
10 & is\_permanent & is\_major\_city & -0.207421 \\
11 & role\_encoded & missing\_province & 0.206943 \\
12 & cpi & missing\_province & -0.162138 \\
13 & is\_senior & missing\_province & 0.142926 \\
14 & role\_encoded & cpi & 0.140449 \\
\end{longtable}

\begin{longtable}[]{@{}llll@{}}
\toprule\noalign{}
& feat\_a & feat\_b & pearson\_r \\
\midrule\noalign{}
\endhead
\bottomrule\noalign{}
\endlastfoot
0 & year & month & -0.918519 \\
\end{longtable}

\subsubsection{Random Forest model}\label{random-forest-model}

\begin{verbatim}
Random Forest CV RMSE (original scale): 39289.29
\end{verbatim}

\begin{verbatim}
Random Forest Test RMSE: 42782.07
Random Forest Test MAE:  27155.26
Random Forest Test R²:   0.3519
\end{verbatim}

\subsubsection{Baseline (mean on log
scale)}\label{baseline-mean-on-log-scale}

\begin{verbatim}
Baseline Test RMSE: 55915.85
Baseline Test MAE:  43253.41
Baseline Test R²:   -0.1070
\end{verbatim}

\subsubsection{Generalized additive model
(GAM)}\label{generalized-additive-model-gam}

A \textbf{pyGAM} \texttt{LinearGAM} uses \textbf{smooth spline terms} on
continuous predictors (salary-based role encoding, CPI, demand, calendar
year and month) and \textbf{linear} terms on the binary indicators. It
is fit on \texttt{log1p(salary\_mid)}; RMSE and related metrics are on
the \textbf{original salary scale}, consistent with Random Forest and
XGBoost. Non-finite values in \texttt{X} are coerced to missing and
imputed with the training median (per CV fold, then once on full
training for the final fit), because pyGAM requires finite numeric
inputs. (Install: \texttt{pip\ install\ pygam}.)

\begin{verbatim}
GAM CV RMSE (original scale): 39713.06
GAM Test RMSE: 39999.67
GAM Test MAE:  25906.58
GAM Test R²:   0.4335
\end{verbatim}

\subsubsection{XGBoost model}\label{xgboost-model}

\begin{verbatim}
Best parameters: {'learning_rate': 0.025, 'max_depth': 4, 'n_estimators': 200}
Best CV RMSE (original scale): 37900.33
Test RMSE: 39749.3604
Test MAE:  25104.5959
Test R²:   0.4406
\end{verbatim}

Importance of Features

\pandocbounded{\includegraphics[keepaspectratio]{midterm_files/figure-pdf/cell-27-output-1.pdf}}

\begin{verbatim}
role_encoded        0.515994
role_demand         0.159155
is_senior           0.095154
year                0.050224
missing_province    0.041479
is_permanent        0.040849
month               0.040351
cpi                 0.030086
is_major_city       0.026708
dtype: float32
\end{verbatim}

Linear Regression Model for dummy Roles

\begin{verbatim}
role_Data Entry       -0.913349
role_Other            -0.129810
role_demand           -0.108063
role_Developer         0.091853
role_Data Scientist    0.154483
role_Data Engineer     0.246001
role_Manager           0.257627
dtype: float64
\end{verbatim}

\pandocbounded{\includegraphics[keepaspectratio]{midterm_files/figure-pdf/cell-30-output-1.pdf}}

Partial Dependence Plots

\pandocbounded{\includegraphics[keepaspectratio]{midterm_files/figure-pdf/cell-31-output-1.pdf}}

\pandocbounded{\includegraphics[keepaspectratio]{midterm_files/figure-pdf/cell-32-output-1.pdf}}

\begin{verbatim}
Unable to display output for mime type(s): text/html
\end{verbatim}

\begin{verbatim}
Unable to display output for mime type(s): text/html
\end{verbatim}

\begin{verbatim}
Unable to display output for mime type(s): text/html
\end{verbatim}

\begin{verbatim}
Unable to display output for mime type(s): text/html
\end{verbatim}

\subsubsection{Role × province (interactive
treemap)}\label{role-province-interactive-treemap}

Tile \textbf{area} is the number of postings in each role--province
cell; \textbf{color} is the \textbf{median} advertised salary in that
cell. Zoom by clicking into a role; hover for exact counts and median
pay.

\begin{verbatim}
Unable to display output for mime type(s): text/html
\end{verbatim}




\end{document}
